\section{Contexto operacional de missões de rendezvous}

As missões de rendezvous em órbita ocupam posição relevante no desenvolvimento de tecnologias espaciais, tanto por sua importância histórica quanto por sua aplicação em operações contemporâneas de suporte, inspeção, manutenção, reabastecimento e serviço em órbita. Ao longo das últimas décadas, essas missões deixaram de ser tratadas apenas como demonstrações pontuais de capacidade orbital e passaram a integrar arquiteturas operacionais mais amplas, nas quais a aproximação segura entre veículos espaciais constitui etapa essencial para a execução da missão \cite{arantes2010far,technicalreportsoyuz}.

No contexto de missões tripuladas, o relatório técnico sobre a espaçonave Soyuz mostra que as operações de rendezvous envolvem diferentes escalas de problema, desde a inserção orbital e os ajustes de transferência até a navegação de proximidade e a fase final de encontro com a Estação Espacial Internacional. O estudo destaca a evolução de diferentes perfis operacionais, incluindo esquemas clássicos, rápidos e ultrarrápidos de rendezvous, bem como o uso das equações de Hill--Clohessy--Wiltshire nas operações de proximidade, evidenciando a busca por maior eficiência temporal, economia de recursos e robustez operacional \cite{technicalreportsoyuz}.

Além do contexto das missões tripuladas, a literatura sobre satélites de serviço mostra que o rendezvous também se tornou elemento central em missões não tripuladas voltadas à inspeção, observação, manutenção e atendimento de espaçonaves-alvo. No trabalho de \cite{arantes2010far}, por exemplo, são analisadas manobras far e proximity de uma constelação de satélites de serviço, combinando transferência orbital, dinâmica relativa e estimação autônoma da pose de um alvo não cooperativo por visão computacional. Essa abordagem ilustra uma mudança importante no contexto operacional do rendezvous, que deixa de se restringir ao encontro entre veículos cooperativos e passa a incluir cenários de serviço em órbita com maior grau de autonomia e maior complexidade operacional \cite{arantes2010far}.

Esse alargamento do contexto operacional também traz novos desafios de guiamento, controle e segurança durante a aproximação. Em operações modernas, não basta conduzir a espaçonave em direção ao alvo; é necessário respeitar restrições geométricas, corredores seguros, limitações dinâmicas e, em alguns casos, a presença de obstáculos na região de aproximação. O trabalho de \cite{fiori2024modeling} ilustra essa perspectiva ao tratar o rendezvous automatizado nas proximidades de uma estação espacial sob restrições posicionais e com obstáculos, considerando ainda fases distintas de aproximação no referencial \gls{lvlh}. Embora esse estudo não tenha foco em manipuladores robóticos, ele mostra que o contexto contemporâneo de rendezvous exige soluções integradas de modelagem, guiamento e controle em cenários mais restritivos \cite{fiori2024modeling}.

De modo geral, observa-se que o contexto operacional das missões de rendezvous evoluiu de operações predominantemente cooperativas e rigidamente planejadas para missões cada vez mais autônomas, flexíveis e orientadas a serviço em órbita. Essa evolução amplia a relevância de estudos voltados à dinâmica relativa, ao guiamento em proximidade, à percepção do alvo e à interação entre espaçonave e manipulador robótico. Para a presente dissertação, esse contexto é importante porque evidencia que a modelagem das fases de aproximação curta e interação final deve ser compreendida como parte de um problema operacional mais amplo, no qual precisão, segurança e autonomia se tornam requisitos centrais.